% TOP_PROBLEM: {"id": 1511, "title": "Euler-Mascheroni Constant Irrationality", "category_id": 1, "category": {"id": 1, "name": "number_theory", "display_name": "Number Theory", "description": "Properties of integers, prime numbers, Diophantine equations.", "slug": "number-theory", "order_index": 1, "created_at": "2026-07-31T15:26:25.670Z"}, "status": "open"}
% FIRST_POSTED: null
% ENTRY_KIND: statement_only
% Imported from: batch08.tex; UnsolvedMath ID 1511
\documentclass[11pt]{article}
\usepackage[margin=25mm]{geometry}
\usepackage{fontspec}
\setmainfont{Latin Modern Roman}
\usepackage{amsmath,amssymb,mathtools}
\usepackage{unicode-math}
\setmathfont{Latin Modern Math}
\usepackage{xurl,hyperref}
\hypersetup{colorlinks=true,urlcolor=blue}
\setcounter{secnumdepth}{0}
\setlength{\parindent}{0pt}
\setlength{\parskip}{5pt}
\setlength{\emergencystretch}{3em}
\newcommand{\sref}[2]{\hyperlink{source:#1:#2}{[S#2]}}
\newcommand{\cataloguescope}[1]{\paragraph{Recorded target.}#1}
\begin{document}
% UnsolvedMath ID 1511; source catalogue code NUM-017
\setcounter{section}{1510}
\section{Euler-Mascheroni Constant Irrationality}\label{1511}
{\small\sffamily UnsolvedMath ID 1511 · Number Theory}
\subsection{Definitions and mathematical statement}

\paragraph{Shared notation.}$\mathbb N=\{1,2,\ldots\}$, and $\mathbb Z,\mathbb Q,\mathbb R,\mathbb C$ denote the integers, rationals, reals and complex numbers. Any other conventions are specified locally.


Let \(H_N=\sum_{j=1}^N1/j\) and let \(\log\) denote
the natural logarithm. The Euler--Mascheroni constant is the real limit
\[
                \gamma=\lim_{N\to\infty}(H_N-\log N).
\]
A real number is irrational if it is not in \(\mathbb Q\).
It is transcendental if it is not a root of any nonzero polynomial
in \(\mathbb Z[X]\); transcendence implies irrationality.
\paragraph{Statement recorded in the supplied source.}
Determine both of the following assertions:
\[
 \gamma\notin\mathbb Q;
 \qquad P(\gamma)\ne0\quad\text{for every nonzero }P\in\mathbb Z[X].
\]
Thus, is \(\gamma\) irrational, and more strongly, is it transcendental? \sref{1511}{2} \sref{1511}{3}
\subsection{Short English statement}
Is the Euler--Mascheroni constant irrational? Is it also transcendental?
\subsection{Sources}
\begin{description}
\item[\hypertarget{source:1511:1}{S1}] ULAM.AI and the UnsolvedMath Contributors, \emph{UnsolvedMath: A Curated Collection of Open Mathematics Problems}, record 1511 (NUM-017). CC BY 4.0. \url{https://huggingface.co/datasets/ulamai/UnsolvedMath}.
\item[\hypertarget{source:1511:2}{S2}] Jonathan Sondow, \emph{Criteria for Irrationality of Euler's Constant}, Proceedings of the American Mathematical Society 131 (2003), 3335--3344; arXiv:math/0209070v2, Section 1, equation (1). \url{https://arxiv.org/abs/math/0209070}.
\item[\hypertarget{source:1511:3}{S3}] Michel Waldschmidt, \emph{Open Diophantine Problems}, arXiv:math/0312440v2 (2004), Section 2.3, the Euler-constant irrationality question, and Section 3.1, the stronger non-period/transcendence question. \url{https://arxiv.org/abs/math/0312440}.
\end{description}
\label{end:1511}
\end{document}
